% matula-213a.tex -- a report on the answer to exercise 7.2.2.3-213(a).
% Typeset with:  tex matula-213a; dvips matula-213a; ps2pdf matula-213a.ps
% The figure matula.ST is made by:  mpost st.mp   (which inputs matula.mp).

\input epsf
\parskip=0pt plus 1pt
\font\eighttt=cmtt8 \font\sixtt=cmtt8 at 6pt
\scriptfont\ttfam=\eighttt \scriptscriptfont\ttfam=\sixtt
\def\arc#1#2{{{\tt #1}\atop{\tt #2}}}
\def\n#1{{\tt #1}}
\def\dash{\thinspace{\vrule height 2.6pt depth -2.2pt width 5pt}\thinspace}
\def\edge#1#2{{\tt #1}\dash{\tt #2}}

\centerline{\bf A slip in the answer to exercise 7.2.2.3--213(a)}
\medskip
\centerline{Pre-Fascicle 7A, the printing of December 5, 2024}
\medskip
\centerline{reported by Soojin Nam}
\bigskip

\noindent
The picture that answers exercise 213(a), on page 205, exhibits a map from $S$
into $T$ that is not an embedding: two of its eighteen edges are not edges
of~$T$, and the picture draws those two nonexistent edges. The claim of
uniqueness just above it needs amending as well. Everything else in
exercise~213 is perfect, including the whole {\tt sol} matrix on page~206;
an account of what was checked appears at the end.

\beginsection The data

Read off the pictures on page 134, node by node, the two trees are these
parent-pointer strings, in the notation that MATULA expects:
$$\vbox{\halign{$#$\hfil&&\quad\hbox{#}\hfil\cr
S&\tt .0111444759a488cfch&($m=19$, $s=4$)\cr
\noalign{\smallskip}
T&\tt .011345676965cc5ffh5cklfn55qjstuuwxxwwuCCuFCpppqrtGOHJRLMNO&($n=59$)\cr}}$$
Nodes of both trees are written in typewriter type below; the context tells
which tree is meant.

\beginsection What the picture on page 205 says

Its labels sit on the following nodes of~$T$:
$$\vbox{\halign{\hbox to 1.05em{\hfil\tt #}&&\hbox to 1.05em{\hfil\tt #}\cr
0&1&2&3&4&5&6&7&8&9&a&b&c&d&e&f&g&h&i\cr
\noalign{\smallskip\hrule\smallskip}
D&C&E&H&u&t&v&w&x&s&j&5&F&y&z&G&Q&O&W\cr}}$$
Sixteen of the eighteen edges of $S$ are honored by this map. The other two are
not: edge \edge ch\ of $S$ goes to \edge FO, and edge \edge fg\ goes to
\edge GQ, whereas neither \edge FO\ nor \edge GQ\ is an edge of~$T$. The
picture nevertheless draws them: it shows 60 distinct edges, although $T$ has
but 58. The true edges \edge GO\ and \edge HQ\ lie in gray right
beside them.

\beginsection Why no local repair is possible

The two branches that node \n4 sends into $T$ have been interchanged. Consider
the subtree $T_e$ for $e=\arc uF$: its nodes are \n F, \n G, \n O, \n P, \n W,
with edges \edge FG, \edge GO, \edge OP, \edge OW. Within it \n F has the
single neighbor \n G, so no node with two children can be mapped to \n F,
and \n c has two. Relabeling cannot help either, since the subtree of $S$
rooted at \n c is the path \edge gf\dash{\tt c}\dash{\tt h}\dash{\tt i} on five
nodes, while the longest path in $T_{\arc uF}$ has only four.

The {\tt sol} matrix on page 206 says precisely this. In its column headed
$\arc uF$, row \n c holds 0 --- and row \n5 holds 1.

\beginsection The correct answer

The path \edge59\dash{\tt a}\dash{\tt b} belongs in the \n F branch, and the
five-node subtree at \n c belongs in the \n t branch:
$$\vbox{\halign{\hbox to 1.05em{\hfil\tt #}&&\hbox to 1.05em{\hfil\tt #}\cr
0&1&2&3&4&5&6&7&8&9&a&b&c&d&e&f&g&h&i\cr
\noalign{\smallskip\hrule\smallskip}
D&C&E&H&u&F&v&w&x&G&O&W&t&y&z&N&V&s&j\cr}}$$
This is what MATULA itself prints, in agreement with its {\tt sol} matrix.
(We ran a line-for-line translation of {\tt matula.w} into Go; it reproduces
the original's mem counts to the digit.)
\predisplaypenalty=0 \postdisplaypenalty=10000
$$\epsfbox{matula.ST}$$
\predisplaypenalty=10000 \postdisplaypenalty=0
\centerline{\it The corrected answer to {\rm 213(a)}; compare the picture on
page 205.}

\beginsection Uniqueness

The answer states that the solution is unique except for permutation of
$\{\n0,\n2,\n3\}$. In fact $S$ can be embedded in $T$ in exactly 48 ways.
There are four independent freedoms:
\medskip
\item{$\bullet$} the six permutations of the leaves $\{\n0,\n2,\n3\}$ at \n1;
\item{$\bullet$} the transposition of the leaves \n d and \n e at \n8;
\item{$\bullet$} the transposition of the two legs \edge fg\ and \edge hi\ at \n c;
\item{$\bullet$} the image of \n b, which may be \n P or \n W.
\medskip
\noindent
The first three generate the automorphism group of $S$, whose order is
$6\cdot2\cdot2=24$; the fourth changes the set of nodes used. So there are
exactly two essentially different answers, and $48=24\cdot2$ labeled ones.
A possible rewording: {\it the answer is unique except for permutation of\/}
$\{\n0,\n2,\n3\}$, {\it transposition of\/} \n d {\it and\/} \n e, {\it
transposition of the legs\/} \edge fg\ {\it and\/} \edge hi, {\it and the
choice of\/} \n P {\it or\/} \n W {\it as the image of\/} \n b.

\beginsection What was checked, and how

Parts (b) through (f) were verified independently, without assuming the theory
that the exercise develops:
\medskip
\item{(b)} $S_{\tt 7}=\{\n7,\n8,\n d,\n e\}$ and
  $T_{\arc uw}=\{\n w,\n x,\n y,\n z,\n A,\n B\}$, and the stated map is one of
  the exactly two embeddings that send \n7 to \n w.
\item{(c)} The characterization, and the bipartite-matching form of it
  given in (d), were compared against plain exhaustive backtracking, which assumes nothing about
  the independence of branches: all $18\times116=2088$ entries agree for the
  trees of (a), as do 17,914 further entries from 300 random pairs of trees.
\item{(e)} The property of Algorithm 7.5.1H --- that a matching of size $k$
  leaving $g_j$ unmatched exists if and only if $g_j$ is in the queue at
  termination --- held in 18,479 random bipartite instances checked against
  exhaustive enumeration, and in every subproblem arising in (a).
\item{(f)} The sketch was implemented afresh from the text of the answer alone.
  It reproduces the printed {\tt sol} matrix in all 2088 entries, solves 86
  matching problems (against the bound $mn=1121$), and never needs more than
  three boys, the maximum inner degree of $S$ being 4. The printed matrix obeys
  conditions (i) and (ii) on the numbering of the arcs, and its row \n1 holds 1
  in exactly three columns, namely $\arc DC$, $\arc EC$, and $\arc HC$.
\medskip
\noindent
Only the picture on page 205, and the sentence just above it, are in error.

\bye
